% Prose for this counterexample.  Every region below is delimited by an
% environment that tools/build.py extracts; machine facts (ids, uid, dates,
% urls, enums) live in case.json.  See counterexamples/AGENTS.md.

\cxtitle{Macdonald polynomial ratio's Schur-convexity}

\begin{cxcredits}
\posedby{McSwiggen and Sahi \citeyearpar{McSwiggenSahi2026}}
\foundby{GPT-5.6 (Pro)}{2026-05-14}
\formalizedby{S.~Sra}
\auditedby{S.~Sra}
\contributedby{S.~Sra}
\end{cxcredits}

\begin{cxcontext}
For parameters $q,t\in(0,1)$ and a partition $\lambda$ with at most $n$ parts, $P_\lambda(x;q,t)$ denotes the Macdonald polynomial in the monic normalization of Macdonald's book, as used in \citet{McSwiggenSahi2026}; at $q=t$ it specializes to the Schur polynomial $s_\lambda(x)$, which is what the refutation below exploits.  The statement concerns two convexity properties of the map $\lambda\mapsto\Omega_\lambda(x;q,t)$ on partitions: \emph{log-convexity}, and \emph{Schur-convexity}, meaning monotonicity with respect to the dominance (majorization) order, so that $\lambda\succeq\mu$ implies $\Omega_\lambda\ge\Omega_\mu$.  The normalization $\Omega_\lambda$ and the lattice $\mathcal L_n$ of evaluation points are defined inside the quoted statement.
\end{cxcontext}

% ---------------- result: macdonald-schur-convexity ----------------

\begin{cxsource}{macdonald-schur-convexity}
C.~McSwiggen and S.~Sahi \citeyearpar{McSwiggenSahi2026}, Theorem 2.1
\end{cxsource}

\begin{cxstatement}{macdonald-schur-convexity}
For all $q,t\in(0,1)$, $a>0$, and $x\in\mathcal L_n^{q,t,a}$, the function $\lambda\mapsto\Omega_\lambda(x;q,t)$ is log-convex and Schur-convex.

Here, by \S2.1 of the same paper, $\Omega_\lambda(x;q,t)=P_\lambda(x;q,t)/P_\lambda(t^\delta;q,t)$ is the Macdonald polynomial normalized to $1$ at the principal specialization $t^\delta=(t^{n-1},t^{n-2},\ldots,1)$, and the scaled, $t$-shifted, dominant $q$-lattice is
\[\mathcal L_n=\mathcal L_n^{q,t,a}=\left\{(aq^{-\mu_1}t^{n-1},aq^{-\mu_2}t^{n-2},\ldots,aq^{-\mu_n}):(\mu_1\ge\ldots\ge\mu_n)\in\mathbb Z^n\right\}\subset\R^n.\]
\end{cxstatement}

\begin{cxsummary}{macdonald-schur-convexity}
McSwiggen--Sahi Theorem 2.1: Macdonald lattice Schur-convexity \citeyearpar{McSwiggenSahi2026}
\end{cxsummary}

\begin{cxcertificate}{macdonald-schur-convexity}
Rank-two Schur specialization, valid for every $0<r<1$.
\end{cxcertificate}

\begin{cxrefutation}
\begin{theorem}[\statusfalse: Theorem 2.1 of \citet{McSwiggenSahi2026}]\label{thm:macdonald}
The stated Schur-convexity fails already in rank two.
\end{theorem}
\begin{proof}
Set $q=t=r\in(0,1)$, $a=1$, and choose lattice index $\nu=(1,0)$.  The resulting lattice point is $x=(1,1)$.  Let $\lambda=(2,0)$ and $\mu=(1,1)$; then $\lambda\succeq\mu$.  At $q=t$, Macdonald polynomials specialize to Schur polynomials, and direct evaluation gives
\[
\Omega_{(2,0)}(1,1;r,r)=\frac{3}{1+r+r^2},\qquad
\Omega_{(1,1)}(1,1;r,r)=\frac1r.
\]
But $1+r+r^2-3r=(1-r)^2>0$, so $\Omega_{(2,0)}<\Omega_{(1,1)}$ for every $0<r<1$, opposite to Schur-convexity.  At $r=1/2$, the values are $12/7<2$.
\end{proof}

\begin{remark}[Scope]
The specialization $q=t$ is a convenience, not the mechanism.  Appendix~\ref{sec:macdonaldomega} carries the reversal for all $q,t\in(0,1)$, with gap $(1-t)^2(1-qt)/\bigl(t(1+t)(1-qt^2)\bigr)$, and shows it is admissible whenever $t=q^k$ for a positive integer $k$ --- the case above being $k=1$.  It also locates the failure, in a determinant shift that the principal-specialization normalization does not commute with, and shows that the $1^n$-normalized ratio $W_\lambda=P_\lambda(x)/P_\lambda(1^n)$ satisfies in this case exactly the inequality $\Omega_\lambda$ violates.  What is refuted is a claim about the log- and Schur-convexity under the normalization used by~\citet{McSwiggenSahi2026}.
\end{remark}
\end{cxrefutation}


%%% Local Variables:
%%% mode: LaTeX
%%% TeX-master: "../../tex/main"
%%% End:
